% RESUMO: um parágrafo, sem citações nem links. Valores vêm de numeros.tex.
Câmeras residenciais costumam enviar o vídeo continuamente para a nuvem, o que consome banda de subida e expõe imagens da casa. Este trabalho mede se a detecção de pessoas pode ser feita no próprio dispositivo de borda, comparando três detectores leves (YOLOv8n, SSD-MobileNet~V1 e EfficientDet-Lite0) em FP32 e INT8, três resoluções de quadro (320, 480 e 640 pixels) e três taxas de câmera (5, 15 e 30 FPS). A borda foi medida em um contêiner ARM64 nativo com 2 núcleos e a nuvem em um servidor x86 com rede emulada, usando as mesmas 350 imagens do COCO com pessoas. A quantização INT8 reduziu a latência por quadro em todas as \nCombBorda{} combinações, com ganho de \razaoYoloTres{} vezes no YOLOv8n, \razaoSsdTres{} no SSD-MobileNet~V1 e \razaoEffdetTres{} no EfficientDet-Lite0, e perda de no máximo cerca de quatro centésimos no mAP@50 da classe pessoa. Em INT8, quase todas as combinações sustentaram 30 FPS; o YOLOv8n teve a melhor qualidade (mAP@50 de \mapYoloFpSeis{} a 640 pixels). Enviar o vídeo à nuvem exigiria de \bandaTresQuinze{} a \bandaSeisQuinze{}\,Mbit/s de subida a 15 FPS, enquanto a borda não envia vídeo. Os tempos vêm de um processador ARM de servidor, não de um Raspberry Pi, e valem como comparação relativa.
